%!TEX encoding = UTF-8 Unicode
% LaTeX Template for CV/Resume
%-------------------------------------------------------------------------------
% LuaLaTeX compilation needed
%-------------------------------------------------------------------------------
% Huanyu Zhou <spica.h.zhou@gmail.com>
% https://spica-vir.github.io/about/
%-------------------------------------------------------------------------------

% configure

\documentclass[a4paper,11pt]{article}
\usepackage{academicCV_zh}

% contacts

\myname{张三}
\myposition{科研工作者}
\myaddress{知名大学骨干院系 米花市 邮箱 000001}

\mytel{123456789}
\myemail{szhang@email.com}
\mysite{this.url.does.not.exist/} % Personal website, without 'https://', optional

% Some self-defined commands to make CV looks better

% Highlight my name in author list
\newcommand{\MeAsAuthor}{\textbf{\underline{Zhang S}}}
\newcommand{\MeAsCorrespondingAuthor}{\textbf{\underline{Zhang S*}}}

% Add comments (editor's pick, under review, etc.) to a paper entry in the publication list
\newcommand{\PaperComment}[1]{\textbf{\underline{(#1)}}}

\begin{document}

%-------------------------------------------------------------------------------
% Contact Info
%-------------------------------------------------------------------------------

\makeheader % basic information defined above

% academic social media
%
% The following commands are defined in a similar way to '\href'.
% For other social media, commands need to be defined in 'my_CVstyle.sty'
% \myGoogleScholar
% \myORCID
% \myResearchGate
% \myAcademia
% \myGitHub
% \myLinkedIn
% \myTwitter
%
\begin{acmedia}
	\myGoogleScholar{https://scholar.google.com/}{Google Scholar}
	\EntrySeparator
	\myORCID{https://orcid.org/}{ORCID}
	\EntrySeparator
	\myGitHub{https://github.com/}{GitHub}
	\EntrySeparator
	\myLinkedIn{https://www.linkedin.com/}{LinkedIn}
\end{acmedia}


%-------------------------------------------------------------------------------
% Employment History
%-------------------------------------------------------------------------------

\section{工作经历}

% an entry of education / employment
% Usage: \experience{<position>}{<date>}{<organization>}{<location>}{<description(optional)>}

\experience
	{科研工作者}
	{2025年至今}
	{知名大学骨干院系}
	{米花市}
	{\begin{itemize}
		\item[-] 关于这个工作的其他信息
	\end{itemize}}

%-------------------------------------------------------------------------------
% Other Positions and Membership of Professional Bodies
%-------------------------------------------------------------------------------

\section{其他社会任职}

\experience
	{会员 / 荣誉会员}
	{2021年至今}
	{国际知名学会}
	{桃花市}{}

\experience
	{技术顾问}
	{2022年至2023年}
	{跨国技术公司}
	{梨花市}{}

%-------------------------------------------------------------------------------
% Education
%-------------------------------------------------------------------------------

\section{教育背景}

\experience
	{某一级学科 工学博士}
	{2021年 $\sim$ 2025年}
	{知名大学骨干院系}
	{米花市}
	{\begin{itemize}
		\item[-] 导师：李四 教授
	\end{itemize}}

\experience
	{某专业 工学硕士}
	{2020年 $\sim$ 2021年}
	{另一知名大学一流学科}
	{樱花市}
	{\begin{itemize}
		\item[-] 导师：王五 教授
	\end{itemize}}

\experience
	{某专业 理学学士}
	{2016年 $\sim$ 2020年}
	{同样知名大学王牌专业}
	{梅花市}
	{\begin{itemize}
		\item[-] 毕业设计导师：赵六 教授
		\item[-] 成绩：\textbf{100}/100，排名：\textbf{前1\%}
	\end{itemize}}

%-------------------------------------------------------------------------------
% Research interests
%-------------------------------------------------------------------------------
\section{研究方向}

\begin{text}
\begin{itemize}[leftmargin=1em]
	\item 拿帽子的研究方向A;
	\item 挣钞票的研究方向B；
	\item 涨知识的研究方向C。
\end{itemize}
\end{text}

%-------------------------------------------------------------------------------
% Awards and Fellowships
%-------------------------------------------------------------------------------

\section{人才项目与科研获奖情况}

% an entry of fellowship / honor / award etc. 
% Usage: \award{<name>}{<date>}{<event>}{<description(optional)>}

\award
	{某重要人才引进计划}
	{2026年至今}
	{引进机构}
	{关于人才计划的其他信息}

\award
	{最佳报告奖}
	{2024年}
	{国际学术会议}{}

%-------------------------------------------------------------------------------
% Grants
%-------------------------------------------------------------------------------

\section{科研项目}

\subsection{主持项目}

% an entry of research grants
% Usage: \grant{<name>}{<date>}{<source>}{<amount(optional)>}{<description(optional)>}

\grant
	{面向某重大战略需求的基础理论研究}
	{2025年 $\sim$ 2030年}
	{慷慨的资助方 （项目编号：AB123456/X）}
	{￥一大笔钱}{}

\subsection{参与项目}

\grant
	{基于某基础理论的应用研究}
	{2021年 $\sim$ 2024年}
	{一家合作公司（项目编号：CD987654/Y）}
	{￥一大笔钱}
	{\begin{itemize}
		\item[-] 合作者：路人甲 博士（某研究所）、龙套乙 副教授（某大学）
	\end{itemize}}

\subsection{计算资源与差旅资助}

\grant
	{国家超算中心计算资源}
	{2022年至今}
	{米花市超算中心（项目编号：01234）}
	{100万核时}{}

\grant
	{国际顶尖会议差旅资助}
	{2023年}
	{会议组织方}
	{￥一大笔钱}{}

%-------------------------------------------------------------------------------
% List of Publications
%-------------------------------------------------------------------------------

\section{论文发表}

% Publication list environment
% Usage: \begin{paperlist}[ascending=<true/false>]\end{paperlist}

\begin{paperlist}[ascending=false]
	% Preprint entry, must be used within paperlist environment
	% Usage: \preprint{<author>}{<title>}{<journal/server>}{<status>}{<year>}{DOI(optional, without https://doi.org/)}{<other comments(optional)>}
	\preprint
		{\MeAsCorrespondingAuthor, Collaborator A, Collaborator B}
		{Preprint for a ground-breaking research on topic A}
		{arXiv}
		{预印本}
		{2026}{}{}

	% Journal paper entry, must be used within paperlist environment
	% Usage: \journal{<author>}{<title>}{<journal>}{<year>}{<volume(optional)>}{<issue(optional)>}{<page(optional)>}{DOI(optional, without https://doi.org/)}{<other comments(optional)>}
	\journal
		{Collaborator A, \MeAsAuthor, Collaborator B}
		{Journal paper for a theoretical research on topic B}
		{Peer-reviewed Journal}
		{2026}
		{00}
		{00}
		{0001--0009}{}{\PaperComment{Cover Article}}

	% Conference paper entry, must be used within paperlist environment
	% Usage: \conference{<author>}{<title>}{<conference>}{<location>}{<date>}{<URL(optional)>}{<other comments(optional)>}
	\conference
		{\MeAsCorrespondingAuthor, Collaborator A}
		{Conference paper for an experimental research on topic C}
		{International Conference on Experimental Research}
		{City, State}
		{11th Dec. 2025}{}{}

	% Book chapter entry, must be used within paperlist environment
	% Usage: \chapter{<author>}{<chaptertitle>}{<editor(optional)>}{<booktitle>}{<edition(optional)>}{<location(optional)>}{<publisher>}{<year>}{<page>}{<URL(optional)>}{<other comments(optional)>}
	\chapter
		{\MeAsAuthor}
		{Fundamental Theory on Topic D}
		{Mr Editor, Ms Editor}
		{Introduction to Topic D}
		{1st}
		{Location}
		{Publisher}
		{2020}
		{16-32}{}{}

\end{paperlist}

%-------------------------------------------------------------------------------
% List of Presentations
%-------------------------------------------------------------------------------

\section{会议报告}

\subsection{特邀报告}
% Presentation list environment
% Usage: \begin{prelist}\end{prelist}
\begin{prelist}
	% keynote presentation entry, must be used within prelist environment
	% Usage: \keynote{<title>}{<conference>}{<location>}{<date>}{<URL(optional)>}{<other comments(optional)>}
	\keynote
		{The progress review of research topic A}
		{International Conference on Topic A}
		{City, State}
		{2024}{https://www.youtube.com/}{}

	% guest presentation entry, must be used within prelist environment
	% Usage: \guest{<title>}{<conference>}{<location>}{<date>}{<URL(optional)>}{<other comments(optional)>}
	\guest
		{My research on topic B}
		{Departmental Seminar}
		{Department, University}
		{2023}{}{}
\end{prelist}

\subsection{投稿报告}
\begin{prelist}
	% oral presentation entry, must be used within prelist environment
	% Usage: \oral{<title>}{<conference>}{<location>}{<date>}{<URL(optional)>}{<other comments(optional)>}
	\oral
		{The presentation of some preliminary results}
		{Regional Conference on Topic C}
		{City, State}
		{2021}{}{}

	% poster presentation entry, must be used within prelist environment
	% Usage: \poster{<title>}{<conference>}{<location>}{<date>}{<URL(optional)>}{<other comments(optional)>}
	\poster
		{To exchange ideas on topic D}
		{Discussion meeting by a society}
		{Academic Society}
		{2022}{}{}

\end{prelist}

%-------------------------------------------------------------------------------
% Teaching and Mentoring experiences
%-------------------------------------------------------------------------------
\section{指导科研与授课}

\subsection{承担课程}

\experience
	{硕士生科研课程}
	{2025年至今}
	{授课教师 知名大学}
	{}{}

\experience
	{本科生实验课}
	{2021年 $\sim$ 2025年}
	{助教 另一知名大学}
	{}{}

\subsection{指导科研项目}

\experience
	{一项有趣的研究}
	{2026年至今}
	{博士 知名大学}
	{小帅同学}
	{}

\experience
	{一项更有趣的研究}
	{2024年 $\sim$ 2026年}
	{硕士 知名大学}
	{小美同学}
	{}

\end{document}
